\section{Lecture 15: Backpropagation and Automatic Differentiation}

\subsection{Motivation}

In the previous tutorial, we manually computed gradients in a small neural network using the chain rule.

\medskip

\noindent
However, modern neural networks may contain:
\begin{itemize}
    \item millions or billions of parameters,
    \item deeply nested compositions of functions.
\end{itemize}

\medskip

\noindent
\textbf{Key question.}  
How can we compute gradients efficiently in such systems?

\medskip

\noindent
\textbf{Answer.}  
Backpropagation provides a systematic and efficient method for computing gradients.

\subsection{Computational Graphs}

A computational graph represents a function as a composition of elementary operations.

\medskip

\noindent
\textbf{Example.}
\[
z = f(x,y) = (x \cdot y + y)^2.
\]

\medskip

\noindent
We can represent this as a graph:
\begin{itemize}
    \item $a = x \cdot y$
    \item $b = a + y$
    \item $z = b^2$
\end{itemize}

\medskip

\noindent
\textbf{Structure.}
\begin{itemize}
    \item Nodes: intermediate variables
    \item Edges: functional dependencies
\end{itemize}

\medskip

\noindent
\textbf{Insight.}  
Complex functions can be decomposed into simple operations.

\subsection{Chain Rule and Reverse-Mode Differentiation}

\textbf{Chain rule.}  
For composed functions:
\[
z = f(g(h(x))),
\]
we have:
\[
\frac{dz}{dx} = \frac{dz}{df} \cdot \frac{df}{dg} \cdot \frac{dg}{dh} \cdot \frac{dh}{dx}.
\]

\medskip

\noindent
\textbf{Forward vs reverse mode.}

\begin{itemize}
    \item Forward mode:
    \begin{itemize}
        \item propagate derivatives from inputs to outputs
    \end{itemize}

    \item Reverse mode (backpropagation):
    \begin{itemize}
        \item propagate gradients from outputs to inputs
    \end{itemize}
\end{itemize}

\medskip

\noindent
\textbf{Key observation.}
\begin{itemize}
    \item Neural networks typically have many parameters and one output
    \item Reverse mode is more efficient in this setting
\end{itemize}

\medskip

\noindent
\textbf{Insight.}  
Backpropagation is reverse-mode automatic differentiation.

\subsection{Efficient Gradient Computation}

Let $L$ be the loss function and $\theta$ the parameters.

\medskip

\noindent
\textbf{Goal.}
\[
\nabla_\theta L.
\]

\medskip

\noindent
\textbf{Backpropagation procedure.}

\begin{enumerate}
    \item \textbf{Forward pass:}
    \begin{itemize}
        \item Compute all intermediate variables
    \end{itemize}

    \item \textbf{Backward pass:}
    \begin{itemize}
        \item Compute gradients starting from output
        \item Use chain rule to propagate backward
    \end{itemize}
\end{enumerate}

\medskip

\noindent
\textbf{Local gradients.}
\begin{itemize}
    \item Each node computes partial derivatives with respect to its inputs
\end{itemize}

\medskip

\noindent
\textbf{Gradient accumulation.}
\[
\frac{\partial L}{\partial x} = \sum_{\text{paths}} \prod \text{local derivatives}.
\]

\medskip

\noindent
\textbf{Efficiency.}
\begin{itemize}
    \item Each intermediate result is reused
    \item Computational cost comparable to forward pass
\end{itemize}

\medskip

\noindent
\textbf{Insight.}  
Backpropagation avoids redundant computation through dynamic programming.

\subsection{Backpropagation in Neural Networks}

Consider a layer:
\[
z_\ell = W_\ell h_{\ell-1} + b_\ell, \quad h_\ell = \sigma(z_\ell).
\]

\medskip

\noindent
\textbf{Backward pass.}

Let $\delta_\ell = \frac{\partial L}{\partial z_\ell}$.

\medskip

\noindent
\textbf{Recursion.}
\[
\delta_\ell = (W_{\ell+1}^\top \delta_{\ell+1}) \odot \sigma'(z_\ell).
\]

\medskip

\noindent
\textbf{Gradients.}
\[
\frac{\partial L}{\partial W_\ell} = \delta_\ell h_{\ell-1}^\top,
\]
\[
\frac{\partial L}{\partial b_\ell} = \delta_\ell.
\]

\medskip

\noindent
\textbf{Interpretation.}
\begin{itemize}
    \item Gradients flow backward through layers
    \item Each layer transforms gradients using local derivatives
\end{itemize}

\medskip

\noindent
\textbf{Insight.}  
Backpropagation is a structured application of the chain rule across layers.

\subsection{Implementation Perspectives}

Modern machine learning frameworks implement automatic differentiation.

\medskip

\noindent
\textbf{Key ideas.}

\begin{itemize}
    \item Build computational graph during forward pass
    \item Store intermediate values
    \item Automatically compute gradients in backward pass
\end{itemize}

\medskip

\noindent
\textbf{Advantages.}

\begin{itemize}
    \item No need to manually derive gradients
    \item Works for arbitrary compositions of functions
\end{itemize}

\medskip

\noindent
\textbf{Tradeoffs.}

\begin{itemize}
    \item Memory cost for storing intermediate values
    \item Numerical stability considerations
\end{itemize}

\medskip

\noindent
\textbf{Insight.}  
Automatic differentiation turns gradient computation into a programmable operation.

\subsection{A Unifying Perspective}

Backpropagation unifies several ideas:

\begin{itemize}
    \item \textbf{Chain rule:} foundation of gradient computation
    \item \textbf{Computation graphs:} structure of functions
    \item \textbf{Dynamic programming:} efficient reuse of computations
    \item \textbf{Optimization:} enables gradient-based learning
\end{itemize}

\medskip

\noindent
\textbf{Core idea.}  
Gradients of complex functions can be computed efficiently by propagating local derivatives backward.

\medskip

\noindent
\textbf{Key insight.}  
Backpropagation makes training deep neural networks computationally feasible.

\subsection{Outlook}

In this lecture, we introduced:

\begin{itemize}
    \item computational graphs,
    \item reverse-mode differentiation,
    \item efficient gradient computation,
    \item implementation via automatic differentiation.
\end{itemize}

\medskip

\noindent
We now turn to the practical challenges of training deep networks.

\medskip

\noindent
In the next lecture, we study initialization, vanishing and exploding gradients, and optimization difficulties in deep learning.

\medskip

\noindent
\textbf{Guiding principle.}  
Efficient gradient computation is the foundation of modern deep learning.